% 数学ノート集・共通テンプレート v2（積分論2026年度第1部の版面を参考）
% A5判。章内で定義・定理・例・注意等の番号を共通化。
% コンパイル： platex -interaction=nonstopmode math_note_template_v2.tex (2回)
%              dvipdfmx math_note_template_v2.dvi
\documentclass[a5paper,10pt,onecolumn]{jsarticle}
\usepackage[dvipdfmx,top=22mm,bottom=20mm,inner=19mm,outer=18mm,headheight=13pt,headsep=6mm]{geometry}
\usepackage{amsmath,amssymb,amsthm,mathtools}
\usepackage{enumitem,fancyhdr,titlesec}
\usepackage[dvipdfmx,bookmarks=true,colorlinks=false]{hyperref}
\usepackage{pxjahyper}

% ----- 差し替える文書情報 -----
\newcommand{\NoteTitle}{de Rhamコホモロジーと大域的位相}
\newcommand{\NoteSeries}{数学ノート集}

% ----- 文字・組版 -----
% 参考PDFの質素な明朝＋Computer Modernの印象を引き継ぎながら，
% A5判で窮屈にならないよう本文の行間を少し確保する．
\linespread{1.13}
\setlength{\parindent}{1zw}
\setlength{\parskip}{0pt}
\setlength{\abovedisplayskip}{6pt plus 2pt minus 2pt}
\setlength{\belowdisplayskip}{6pt plus 2pt minus 2pt}
\setlength{\abovedisplayshortskip}{3pt plus 2pt minus 1pt}
\setlength{\belowdisplayshortskip}{4pt plus 2pt minus 1pt}
\setlist[enumerate]{topsep=3pt,itemsep=1pt,parsep=0pt,leftmargin=2.4zw}
\titleformat{\section}{\normalfont\normalsize\bfseries}{\thesection}{0.8zw}{}
\titlespacing*{\section}{0pt}{2.0ex plus .4ex}{.9ex}
\titleformat{\subsection}{\normalfont\normalsize\bfseries}{\thesubsection}{0.7zw}{}
\titlespacing*{\subsection}{0pt}{1.5ex plus .3ex}{.6ex}
\raggedbottom

% ----- 柱とページ番号：ページ上部、罫線なし -----
\fancypagestyle{mathnote}{%
  \fancyhf{}%
  \fancyhead[L]{\footnotesize \NoteSeries\ （\NoteTitle）}%
  \fancyhead[R]{\footnotesize \thepage}%
  \renewcommand{\headrulewidth}{0pt}%
  \renewcommand{\footrulewidth}{0pt}%
}
\pagestyle{mathnote}

% ----- 節内通し番号の定義・定理群 -----
\newtheoremstyle{mathnotetheorem}
 {4.5pt}{4.5pt}{\normalfont}{}%
 {\bfseries}{}{0.65zw}%
 {\thmname{#1}\thmnumber{ #2}\thmnote{（#3）}}
\theoremstyle{mathnotetheorem}
\newtheorem{theorem}{定理}[section]
\newtheorem{definition}[theorem]{定義}
\newtheorem{lemma}[theorem]{補題}
\newtheorem{proposition}[theorem]{命題}
\newtheorem{corollary}[theorem]{系}
\newtheorem{example}[theorem]{例}
\newtheorem{remark}[theorem]{注意}
\renewcommand{\proofname}{証明}
\renewcommand{\qedsymbol}{$\square$}
\makeatletter
\renewenvironment{proof}[1][\proofname]{%
  \par\pushQED{\qed}\normalfont\topsep4\p@\@plus2\p@\relax
  \trivlist\item[\hskip\labelsep\bfseries #1]\ignorespaces
}{\popQED\endtrivlist\@endpefalse}
\makeatother

% ----- 使うノートで適宜拡張する記号 -----
\newcommand{\R}{\mathbb{R}}
\newcommand{\norm}[1]{\lVert #1\rVert}
\newcommand{\abs}[1]{\lvert #1\rvert}

\begin{document}
\section*{読了後にできること}
本稿を読むための前提は，de Rhamコホモロジー，閉形式と完全形式．
本稿の具体的な到達課題は，閉形式が大域的に完全でないことを積分で証明する．

\section{de Rham複体}

本稿では $M$ を第二可算Hausdorffな滑らかな多様体とし，
微分形式，外微分，単位の分割，
Poincar\'eの補題およびStokesの定理を前提とする．
実係数のコホモロジーを扱う．

\begin{definition}
$k$-次閉形式全体と完全形式全体を
\[
 Z^k(M)=\ker(d:\Omega^k(M)\to\Omega^{k+1}(M)),
\]
\[
 B^k(M)=\operatorname{im}
 (d:\Omega^{k-1}(M)\to\Omega^k(M))
\]
と置く．$d^2=0$ により $B^k(M)\subset Z^k(M)$ なので，
\[
 H_{\mathrm{dR}}^k(M)=Z^k(M)/B^k(M)
\]
を第 $k$ de Rhamコホモロジー群という．
$B^0(M)=0$ と約束する．
\end{definition}

\begin{proposition}[次数零]\label{prop:h0}
$H_{\mathrm{dR}}^0(M)$ は
$M$ の各連結成分上で定数となる滑らかな関数全体である．
特に $M$ が連結なら
$H_{\mathrm{dR}}^0(M)\cong\mathbb R$ である．
\end{proposition}
\begin{proof}
$0$-形式は滑らかな関数である．
$df=0$ なら任意の滑らかな曲線 $\gamma$ に対して
$(f\circ\gamma)'=df_{\gamma(t)}(\gamma'(t))=0$
となるので，$f$ は曲線の像の上で定数である．
多様体は局所道連結なので，
連結成分は道連結成分に等しい．
従って $df=0$ なら各連結成分で定数である．
逆に各連結成分で定数の関数は
連結成分が開であることから滑らかであり，
その外微分は零である．
$B^0=0$ だから結論が従う．
\end{proof}

\begin{proposition}[引き戻しとコホモロジー]\label{prop:pullback}
滑らかな写像 $F:M\to N$ は
\[
 F^*:H_{\mathrm{dR}}^k(N)\to H_{\mathrm{dR}}^k(M)
\]
を誘導し，
$(G\circ F)^*=F^*\circ G^*$ を満たす．
\end{proposition}
\begin{proof}
$d(F^*\omega)=F^*(d\omega)$ が成り立つので，
閉形式は閉形式へ，完全形式は完全形式へ写る．
従って商空間の写像がwell-definedとなる．
微分形式の引き戻しの函手性から
合成についての等式も従う．
\end{proof}

\section{ホモトピー不変性}

\begin{theorem}[ホモトピー公式]\label{thm:homotopy}
$F_0,F_1:M\to N$ を滑らかな写像とし，
$H:[0,1]\times M\to N$ を滑らかなホモトピーとする．
このとき線形写像
$K:\Omega^k(N)\to\Omega^{k-1}(M)$ が存在して
\[
 F_1^*-F_0^*=dK+Kd
\]
を満たす．
\end{theorem}
\begin{proof}
$H^*\omega$ を $[0,1]\times M$ 上の形式として考える．
局所座標において一意に
\[
 H^*\omega=dt\wedge\alpha_t+\beta_t
\]
と分解できる．
ここで $\alpha_t$ は $M$ 上の $(k-1)$-形式，
$\beta_t$ は $k$-形式であり，
$t$ に滑らかに依存する．
\[
 K\omega=\int_0^1\alpha_t\,dt
\]
と定める．
$M$ 上の外微分を $d_M$ とすると
\[
 d_{[0,1]\times M}(dt\wedge\alpha_t+\beta_t)
 =dt\wedge(\partial_t\beta_t-d_M\alpha_t)
 +d_M\beta_t.
\]
引き戻しと外微分の可換性から
$H^*(d\omega)=d(H^*\omega)$ であり，
$K(d\omega)=\int_0^1(\partial_t\beta_t-d_M\alpha_t)\,dt$
となる．従って
\begin{align*}
 d(K\omega)+K(d\omega)
 &=\int_0^1 d_M\alpha_t\,dt
   +\int_0^1(\partial_t\beta_t-d_M\alpha_t)\,dt\\
 &=\beta_1-\beta_0=F_1^*\omega-F_0^*\omega.
\end{align*}
これが主張である．
\end{proof}

\begin{corollary}[de Rhamコホモロジーのホモトピー不変性]
\label{cor:homotopy}
滑らかにホモトピックな写像は
同じde Rhamコホモロジー準同型を誘導する．
特に滑らかなホモトピー同値は
de Rhamコホモロジーの同型を誘導する．
\end{corollary}
\begin{proof}
閉形式 $\omega$ に対して
定理\ref{thm:homotopy}から
$F_1^*\omega-F_0^*\omega=d(K\omega)$
となる．
従って商空間 $H_{\mathrm{dR}}^k$ では
両者の像が一致する．
ホモトピー逆写像が存在する場合には，
誘導写像の合成が恒等となるので同型である．
\end{proof}

\begin{corollary}[収縮可能な多様体]\label{cor:contractible}
$M$ が滑らかに収縮可能なら
\[
 H_{\mathrm{dR}}^0(M)\cong\mathbb R,\qquad
 H_{\mathrm{dR}}^k(M)=0\quad(k\ge1)
\]
となる．
\end{corollary}
\begin{proof}
$M$ と一点空間の間の滑らかなホモトピー同値を取る．
系\ref{cor:homotopy}によってde Rhamコホモロジーは
一点空間のものと同型である．
一点空間では0次の滑らかな関数は実数全体であり，
正の次数の微分形式は零なので結論が従う．
\end{proof}

\section{Mayer--Vietorisの短完全列}

\begin{theorem}[微分形式の短完全列]\label{thm:mv-short}
$M=U\cup V$ を開被覆とする．
各次数 $k$ に対し
\[
 0\to\Omega^k(M)
 \xrightarrow{r}\Omega^k(U)\oplus\Omega^k(V)
 \xrightarrow{s}\Omega^k(U\cap V)\to0
\]
は短完全列である．
ここで
\[
 r(\omega)=(\omega|_U,\omega|_V),\qquad
 s(\alpha,\beta)=\alpha|_{U\cap V}-\beta|_{U\cap V}
\]
と定める．
\end{theorem}
\begin{proof}
$r$ は単射である．
実際 $\omega$ が $U,V$ の両方で零なら
$U\cup V=M$ 上で零である．
$s\circ r=0$ は定義から従う．
$s(\alpha,\beta)=0$ なら
$\alpha,\beta$ は $U\cap V$ で一致するので，
微分形式の層の貼り合わせ公理から
唯一の $\omega\in\Omega^k(M)$ が存在して
$r(\omega)=(\alpha,\beta)$ となる．
従って $\ker s=\operatorname{im}r$ である．

$s$ が全射であることを示す．
$\gamma\in\Omega^k(U\cap V)$ を取り，
開被覆 $\{U,V\}$ に従属する滑らかな単位の分割
$\rho_U+\rho_V=1$ を選ぶ．
$U\cap V$ 上の
$\rho_V\gamma$ は $U$ 上の形式へ零で延長できる．
実際 $\rho_V$ の台は $V$ 内にあるため，
$U\setminus V$ の各点の近傍で
$\rho_V\gamma$ は零となる．
同様に $\rho_U\gamma$ を $V$ 上へ零延長できる．
\[
 \alpha=\rho_V\gamma\text{ の }U\text{ への零延長},\quad
 \beta=-\rho_U\gamma\text{ の }V\text{ への零延長}
\]
と置くと
$s(\alpha,\beta)=(\rho_V+\rho_U)\gamma=\gamma$
となる．従って全射である．
\end{proof}

\begin{corollary}[Mayer--Vietoris長完全列]\label{cor:mv}
$M=U\cup V$ に対して自然な長完全列
\[
 \cdots\to H_{\mathrm{dR}}^k(M)\to
 H_{\mathrm{dR}}^k(U)\oplus H_{\mathrm{dR}}^k(V)\to
 H_{\mathrm{dR}}^k(U\cap V)\xrightarrow{\delta}
 H_{\mathrm{dR}}^{k+1}(M)\to\cdots
\]
が存在する．
\end{corollary}
\begin{proof}
定理\ref{thm:mv-short}の写像 $r,s$ は
制限写像で定義されているので外微分と可換である．
従って微分形式の複体の短完全列となる．
ホモロジー代数における短完全複体の長完全列の定理を
コホモロジーに適用すれば結論が従う．
\end{proof}

\section{球面のde Rhamコホモロジー}

\begin{theorem}[球面の計算]\label{thm:sphere}
$n\ge1$ に対して
\[
 H_{\mathrm{dR}}^k(S^n)\cong
 \begin{cases}
 \mathbb R&(k=0,n),\\
 0&(\text{その他})
 \end{cases}
\]
が成立する．
ただし $n=0$ では $S^0$ は二点集合なので
$H_{\mathrm{dR}}^0(S^0)\cong\mathbb R^2$ である．
\end{theorem}
\begin{proof}
まず $S^1$ を北極と南極を除いた二つの開弧
$U,V$ の合併とする．
$U,V$ は収縮可能であり，
$U\cap V$ は二つの連結成分を持ち，
それぞれ収縮可能である．
従って
\[
 H^0(U)\cong H^0(V)\cong\mathbb R,\qquad
 H^0(U\cap V)\cong\mathbb R^2
\]
であり，正の次数では零となる．
Mayer--Vietoris長完全列は
\[
 0\to H^0(S^1)\to\mathbb R^2
 \xrightarrow{r}\mathbb R^2
 \to H^1(S^1)\to0
\]
を与える．
写像 $r$ は定数関数の差なので
$r(a,b)=(a-b,a-b)$ であり，
その核と余核はそれぞれ $\mathbb R$ と同型である．
従って $H^0(S^1)\cong H^1(S^1)\cong\mathbb R$ である．

$n\ge2$ とする．
球面の北極と南極をそれぞれ除いた
$U,V$ は $\mathbb R^n$ と微分同相なので収縮可能である．
$U\cap V$ は $S^{n-1}$ と滑らかにホモトピー同値である．
Mayer--Vietoris長完全列から，
$k\ge2$ について
\[
 H_{\mathrm{dR}}^k(S^n)
 \cong H_{\mathrm{dR}}^{k-1}(S^{n-1})
\]
となる．
また $U\cap V$ が連結なので，
零次から一次への長完全列を考えると
$H_{\mathrm{dR}}^1(S^n)=0$ となる．
帰納法によって $H_{\mathrm{dR}}^n(S^n)\cong\mathbb R$
かつ $1\le k<n$ のコホモロジーは零となる．
$S^n$ は連結なので $H^0(S^n)\cong\mathbb R$ である．
$n$ より大きい次数の微分形式は零なので
高次のコホモロジーも零である．
\end{proof}

\section{トーラスのコホモロジー}

\begin{theorem}[二次元トーラスの計算]\label{thm:torus}
$T^2=\mathbb R^2/(2\pi\mathbb Z)^2$ とすると，
\[
 H_{\mathrm{dR}}^0(T^2)\cong\mathbb R,\qquad
 H_{\mathrm{dR}}^1(T^2)\cong\mathbb R^2,\qquad
 H_{\mathrm{dR}}^2(T^2)\cong\mathbb R.
\]
一次の基底は $[dx]$，$[dy]$，二次の基底は
$[dx\wedge dy]$ で与えられる．
\end{theorem}
\begin{proof}
$T^2$ 上の滑らかな関数と微分形式の係数は
$2\pi$-周期の滑らかな関数とみなせる．
そのFourier係数は任意の階数で急減少する．
これは部分積分を任意回数行うことで，
各係数が
$O((1+m^2+n^2)^{-N})$ となることから従う．
従ってFourier級数は任意階の導関数について
絶対・一様収束し，項別微分できる．

閉1形式 $\omega=A(x,y)\,dx+B(x,y)\,dy$ を取る．
閉形式の条件は $\partial_xB=\partial_yA$ である．
Fourier係数を $A_{mn},B_{mn}$ と書くと
\[
 mB_{mn}=nA_{mn}
\]
が全ての $(m,n)\in\mathbb Z^2$ について成立する．
$(m,n)\ne(0,0)$ なら
\[
 u_{mn}=
 \frac{mA_{mn}+nB_{mn}}{i(m^2+n^2)}
\]
と置く．
閉形式条件から
$i m u_{mn}=A_{mn}$，
$i n u_{mn}=B_{mn}$ が従う．
係数は急減少するので，
$u=\sum_{(m,n)\ne(0,0)}
 u_{mn}e^{i(mx+ny)}$ は滑らかであり，
\[
 du=(A-A_{00})\,dx+(B-B_{00})\,dy.
\]
従って全ての閉1形式は
$A_{00}dx+B_{00}dy$ と完全形式の和である．
一方，$dx$ と $dy$ の線形結合が完全なら，
それぞれ $x$ 方向・$y$ 方向の基本閉曲線で積分すると
係数が零でなければならない．
よって $[dx],[dy]$ は一次の基底となる．

二次の形式は $f(x,y)\,dx\wedge dy$ の形であり，
自動的に閉形式である．
$f_{mn}$ をそのFourier係数とする．
$(m,n)\ne(0,0)$ について
\[
 \eta_{mn}=
 \frac{i f_{mn}}{m^2+n^2}
 e^{i(mx+ny)}(n\,dx-m\,dy)
\]
と置くと
\[
 d\eta_{mn}=f_{mn}e^{i(mx+ny)}dx\wedge dy.
\]
係数の急減少性から
$\eta=\sum_{(m,n)\ne(0,0)}\eta_{mn}$ は滑らかであり，
$d\eta=(f-f_{00})dx\wedge dy$ となる．
従って任意の二次形式は定数倍の
$dx\wedge dy$ と完全形式の和である．
もし $dx\wedge dy$ が完全なら
Stokesの定理から
$\int_{T^2}dx\wedge dy=0$ となるはずだが，
実際には $(2\pi)^2\ne0$ である．
従って二次コホモロジーの基底を与える．
零次については連結性から $\mathbb R$ である．
\end{proof}

\section{微分形式の層とfine resolution}

\begin{definition}
$M$ の開集合 $U$ に対して
$\boldsymbol\Omega^k(U)=\Omega^k(U)$ と置き，
微分形式の制限を射とする．
これによってAbel群の層 $\boldsymbol\Omega^k$ が定まる．
また $\underline{\mathbb R}$ を局所定数実数値関数の層とする．
\end{definition}

\begin{theorem}[Poincar\'eの補題による層の完全分解]
\label{thm:form-resolution}
層の列
\[
 0\to\underline{\mathbb R}\to
 \boldsymbol\Omega^0\xrightarrow d
 \boldsymbol\Omega^1\xrightarrow d\cdots
 \xrightarrow d\boldsymbol\Omega^n\to0
\]
は完全である．
\end{theorem}
\begin{proof}
層の完全性は各茎で検査できる．
次数0では $df=0$ の滑らかな関数は局所定数なので
$\ker(d:\boldsymbol\Omega^0\to\boldsymbol\Omega^1)
=\underline{\mathbb R}$ である．
任意の $p\in M$ では，十分小さい座標球 $U$ を取れる．
座標球は星状開集合と微分同相であるため，
Poincar\'eの補題により $U$ 上の閉 $k$-形式（$k\ge1$）は
$U$ 上の $(k-1)$-形式の外微分となる．
従って各茎で
$\ker d=\operatorname{im}d$ となり，
層の列は全ての次数で完全である．
\end{proof}

\begin{definition}
Abel群の層 $\mathcal F$ が\emph{fine}であるとは，
任意の局所有限開被覆 $\{U_i\}$ に対し，
層の自己準同型 $\theta_i:\mathcal F\to\mathcal F$ で，
$\theta_i$ が $U_i$ の外で零となり，
$\sum_i\theta_i=\operatorname{id}_{\mathcal F}$
が成立するものが存在することをいう．
\end{definition}

\begin{proposition}[微分形式の層はfine]\label{prop:fine-forms}
任意の $k$ について $\boldsymbol\Omega^k$ はfine層である．
\end{proposition}
\begin{proof}
局所有限開被覆 $\{U_i\}$ に従属する
滑らかな単位の分割 $\{\rho_i\}$ を取る．
$\theta_i(\omega)=\rho_i\omega$ と定めると，
$\theta_i$ は制限と可換な層の準同型であり，
$\operatorname{supp}\theta_i\subset U_i$ となる．
さらに局所有限性から
$\sum_i\theta_i(\omega)
=(\sum_i\rho_i)\omega=\omega$ である．
従ってfine層である．
\end{proof}

\begin{lemma}[fine層の非輪状性]\label{lem:fine-acyclic}
パラコンパクトHausdorff空間 $X$ 上のfineなAbel群の層
$\mathcal F$ に対して
\[
 H^q(X,\mathcal F)=0\qquad(q\ge1).
\]
\end{lemma}
\begin{proof}
局所有限な開被覆 $\{U_i\}$ とfine層の自己準同型
$\theta_i$ を取る．
\v Cechコチェイン $c$ に対し，
\[
 (hc)_{i_0\cdots i_{q-1}}
 =\sum_j\theta_j
 \bigl(c_{j i_0\cdots i_{q-1}}\bigr)
\]
と定める．
右辺の各項は $U_j\cap U_{i_0}\cap\cdots\cap U_{i_{q-1}}$
上で定義されるが，$\theta_j$ の台が $U_j$ 内にあるため
必要に応じて零延長して
$U_{i_0}\cap\cdots\cap U_{i_{q-1}}$ 上の切断にできる．
局所有限性によって和は局所的に有限である．
\v Cech微分の交代和を展開すると，
$\delta h+h\delta=\operatorname{id}$ が
正の次数で成立する．
実際，二つの頂点を除く各項は符号を変えて打ち消し合い，
残る項は $\sum_j\theta_j(c)=c$ となる．
従って局所有限な開被覆での正の次数の
\v Cechコホモロジーは零である．
パラコンパクトHausdorff空間では，
開被覆を局所有限細分することができ，
層の導来函手コホモロジーは開被覆を細分したときの
\v Cechコホモロジーの帰納極限として計算できる．
この比較は，flabby分解の\v Cech二重複体を用いて
横・縦のコホモロジーを順に計算することで得られる．
従って全ての正の次数で
$H^q(X,\mathcal F)=0$ となる．
\end{proof}

\begin{theorem}[de Rham複体は層コホモロジーを計算する]
\label{thm:de-rham-sheaf}
任意の滑らかな多様体 $M$ について
\[
 H_{\mathrm{dR}}^k(M)\cong
 H^k(M,\underline{\mathbb R})
\]
が自然に成り立つ．
\end{theorem}
\begin{proof}
定理\ref{thm:form-resolution}により
\[
 0\to\underline{\mathbb R}\to
 \boldsymbol\Omega^0\to\boldsymbol\Omega^1\to\cdots
\]
は層の完全分解である．
命題\ref{prop:fine-forms}と
補題\ref{lem:fine-acyclic}から，
各 $\boldsymbol\Omega^k$ は
大域切断函手に関して非輪状である．
非輪状分解による導来函手の計算定理によって，
$\underline{\mathbb R}$ の層コホモロジーは
この分解に大域切断を適用した複体
\[
 \Omega^0(M)\xrightarrow d\Omega^1(M)\xrightarrow d\cdots
\]
のコホモロジーと自然同型である．
右辺はde Rhamコホモロジーの定義に他ならない．
\end{proof}

\section{特異コホモロジーとの比較}

\begin{definition}
滑らかな特異 $k$-単体
$\sigma:\Delta^k\to M$ と
$\omega\in\Omega^k(M)$ に対して
\[
 I(\omega)(\sigma)=\int_{\Delta^k}\sigma^*\omega
\]
と定める．
これは滑らかな特異コチェインを与える．
\end{definition}

\begin{proposition}[積分写像はコチェイン写像]\label{prop:integration-chain}
任意の $\omega\in\Omega^k(M)$ に対し
\[
 \delta I(\omega)=I(d\omega)
\]
が成立する．
\end{proposition}
\begin{proof}
任意の滑らかな特異 $(k+1)$-単体
$\sigma:\Delta^{k+1}\to M$ に対して
\begin{align*}
 (\delta I(\omega))(\sigma)
 &=I(\omega)(\partial\sigma)\\
 &=\int_{\partial\Delta^{k+1}}\sigma^*\omega\\
 &=\int_{\Delta^{k+1}}d(\sigma^*\omega)\\
 &=\int_{\Delta^{k+1}}\sigma^*(d\omega)\\
 &=I(d\omega)(\sigma).
\end{align*}
三行目は境界付き単体（角を持つ多様体）に対する
Stokesの公式であり，
通常の境界付き多様体の公式を各面に適用して得られる．
従って任意の単体で評価した値が一致し，
コチェイン写像である．
\end{proof}

\begin{theorem}[de Rhamの定理]\label{thm:de-rham}
任意の滑らかな多様体 $M$ について
積分写像は自然同型
\[
 H_{\mathrm{dR}}^k(M)\xrightarrow{\ \cong\ }
 H_{\mathrm{sing}}^k(M;\mathbb R)
\]
を誘導する．
\end{theorem}
\begin{proof}
滑らかな特異コチェインの前層を
$U\mapsto C^k_{\mathrm{sm}}(U;\mathbb R)$ とし，
その層化を $\mathcal S^k$ とする．
各点の十分小さい座標球は収縮可能であるから，
特異コホモロジーのホモトピー不変性により
局所的なコチェイン複体は正の次数で完全であり，
零次の核は局所定数関数である．
従って
\[
 0\to\underline{\mathbb R}\to
 \mathcal S^0\to\mathcal S^1\to\cdots
\]
は層の完全分解となる．
$\mathcal S^k$ はflabbyである：
開集合上の局所コチェインを，十分細かな単体に対する値を
選んで貼り合わせ，残る単体には零を割り当てることで
大きな開集合へ延長できる．
従って層コホモロジーをこの分解から計算できる．

さらに滑らかな特異コチェインからその層化への自然写像は，
大域コホモロジーで同型を誘導する．
実際，任意の開被覆に対して
特異単体の十分細かい重心細分を取れば，
全ての小単体が被覆のどれかの元に含まれる．
重心細分写像は鎖ホモトピーによって恒等写像と
結ばれるため，局所的なコチェインを貼り合わせる操作は
コホモロジーを変えない．
従って
\[
 H^k(M,\underline{\mathbb R})
 \cong H_{\mathrm{sing}}^k(M;\mathbb R)
\]
が得られる．
一方，定理\ref{thm:de-rham-sheaf}によって
左辺はde Rhamコホモロジーである．

積分写像がこの同型と一致することを確認する．
微分形式の層分解から滑らかな特異コチェインの層分解へ，
各局所形式の単体上の積分によって
コチェイン写像が定まる．
命題\ref{prop:integration-chain}から微分と可換であり，
次数零の局所定数関数に対してはその値を保つ．
従って二つの $\underline{\mathbb R}$ の分解を結ぶ
恒等写像の持ち上げとなり，
導来函手の比較定理から誘導写像は上記の自然同型に等しい．
これで積分写像がコホモロジー同型を与える．
\end{proof}

\begin{remark}
de Rhamの定理は，微分形式の局所的な微分方程式
$d\omega=0$ と，空間の大域的な位相的不変量を結びつける．
Poincar\'eの補題は局所的な完全性を保証するが，
大域的な完全性の障害はde Rhamコホモロジー群として現れる．
\end{remark}

\section{到達課題の完全解答}
\noindent\textbf{原点を除いた平面の一次コホモロジー．}
\begin{proposition}
$U=\mathbb R^2\setminus\{0\}$ 上の1形式
\[
\omega=\frac{-y\,dx+x\,dy}{x^2+y^2}
\]
は閉形式であるが完全形式でない．
\end{proposition}
\begin{proof}
$P=-y/(x^2+y^2)$，$Q=x/(x^2+y^2)$ と置くと
$\partial_xQ=(y^2-x^2)/(x^2+y^2)^2$，
$\partial_yP=(y^2-x^2)/(x^2+y^2)^2$．
従って $d\omega=0$ である．
単位円周 $\gamma(t)=(\cos t,\sin t)$ では
$\gamma^*\omega=dt$ だから
$\int_\gamma\omega=2\pi$．
もし $\omega=df$ なら閉曲線上の積分は
$\int_\gamma df=0$ となるので矛盾する．
\end{proof}
\noindent\textbf{到達確認．}
$\omega$ の定数倍 $3\omega$ は完全形式か．
\par\noindent\textbf{解答．}
単位円周上の積分は $6\pi\ne0$ なので完全でない．
\end{document}